\section{Lagrange}

%{\bf Conjugate values of Multiple-valued Function and Conjugate Groups}\\

\begin{theorem}[Lagrange] The order of a subgroup is a divisor of the order of the group.
\end{theorem}

\begin{proof}
Consider a group $G$ of $N$ and a subgroup $H$ composed of the substitutions
\be h_1 = I, h2, h3, \hdots, h_P. \label{HL1} \ee
If G constains no further substitutions, then $N = P$ and the theorem is true. Now let $G$ contain a
substitution $g_2$ not in $H$. Then G also must contain
\be g_2, h_2 g_2, h_3 g_2, \hdots, h_P g_2. \label{HL2} \ee This set are all distinct and all different from the substitutions
(\ref{HL1}), since $h_\alpha g_2 = h_\beta$ requires that $g_2 = h_\alpha^{-1} h_\beta = $ a substition of $H$
contrary to the assumption that $g_2$ is distinct from the elements of $H$. Hence the substitutions combined now
give $2P$ distinct substitutions of $G$. If there are no other substitutions in $G$, then $N=2P$ and the theorem
is true. If there are any other distrinct substitutions $g_3$ in $G$ not in (\ref{HL1}) or (\ref{HL2}), then $G$ contains
\be g_3, h_2 g_3, h_3 g_3, \hdots, h_P g_3. \label{HL3} \ee Moreover, these substitutions are all different from (\ref{HL1})
and (\ref{HL2}), since $h_\alpha g_3 = h_\beta g_2$ requires that $g_3 = h_\alpha^{-1}h_\beta g_2$ and shall belong
to (\ref{HL2}), contrary to assumption. We now have $3P$ distinct substitutions of $G$. Either $N=3P$ or else $G$
contains a substitution $g_4$ not in the sets (\ref{HL1}), (\ref{HL3}), (\ref{HL3}). In the latter case $G$ contains the products
\be g_4, h_2 g_4, h_2 g_4, \hdots, h_P g_4, \label{HL4} \ee all of which are distrinct and all different from the sets
(\ref{HL1}), (\ref{HL3}), (\ref{HL3}), so that we have $4P$ distinct substitutions. Proceeding in this way, we finally
reach a last set of $P$ substitutions
\be g_nu, h_2 g_nu, h_3 g_nu, \hdots, h_P g_\nu, \label{HLP} \ee since the order of $H$ is finite. Hence $N=\nu P$.
\end{proof}


\begin{definition}
The number $\nu$ is called the {\it index} of the subgroup $H$ under $G$ and is denoted by $\nu = [G \mathbin{:} H]$.
\end{definition}

\begin{corollary}
The order of any group $H$ of substitutions on $n$ letters is a divisor of $n!$. $H$ is a subgroup of the symmetric group $G_n$
on $n$ letters.
\end{corollary}

\begin{theorem}
The period of any substitution contained in a group $G$ of order $N$ is a divisor of $N$.
\end{theorem}

\begin{corollary}[Applied to the Symmetric Group]
The order of any substitution group on a set of $n$ elements is an exact divisor of $N=n!$, the quotient
being the number of distinct value of the corresponding multiple-valued function.
\end{corollary}

Let us arrange the $N$ substitutions of $G$ in a rectangular array using the substitutions of a subgroup $H$
in the first row:

\be
\begin{array}{lllll}
h_1=I & h_2 & h_3 & \hdots & h_P \\
g_2 & h_2 g_2 & h_3 g_2 & \hdots & h_P g_2 \\
g_3 & h_2 g_3 & h_3 g_3 & \hdots & h_P g_3 \\
\vdots & \hdots & \hdots & \hdots & \vdots\\
g_\nu & h_2 g_\nu & h_3 g_\nu & \hdots & h_P g_\nu \\
\end{array}\label{Lagrange}
\ee

The $g_1 = I, g_2, g_3, \hdots, g_\nu$ multiply the elements of $H$ on the right hand side forming right cosets.
Similarly, a rectangular array for the substitiions could be formed by using left costs.

\begin{theorem}
If $\psi$ is a rational funciton of $x_1, x_2, \hdots, x_n$ belonging to a subgroup $H$ of index $\nu$, then $\psi$
is $\nu$-valued under the full set of substitutions of $G$.
\end{theorem}

\begin{proof}
Apply to $\psi$ all the $N$ substitutions of $G$ arranged in a rectangular array as in (\ref{Lagrange}). All the substitutions
belonging to the same row give the same answer since
\[ \psi_{h_i g_\alpha} = (\psi_{h_i} )_{g_\alpha} = (\psi)_{g_\alpha} = \psi_{g_\alpha}. \]
Hence there result at most $\nu$ values. But, if
\[ \psi_{g_\alpha} = \psi_{g_\beta} \]
then $\psi_{g_\alpha g_beta^{-1}} = \psi,$ so that $g_\alpha g_\beta^{-1} h_i \in H$, leaving $\psi$ unaltered. Hence $g_\alpha = h_i g_\alpha$, contrary to the assumption made
in forming the rectangular array.
\end{proof}

Regarding the orders of a group $H$ and a larger group $G$ containing $H$ as a subgroup; that is to say, the order $r$ of $H$ is an exact divisor of the order $G$ which contains $H$ as a subgroup, the quotient $m$ being the number of distinct values of a multiple-valued function unaltered by the substitutions of $G$ which are obtained by the substutions of $H$.

\begin{definition}
The $\nu$ distinct functions $\psi, \psi_{g_2}, \psi_{g_3}, \hdots, \psi_{g_\nu}$ are called {\bf conjugate values} of $\psi$ under $G$.
\end{definition}

\begin{theorem}
The $\rho$ distinct values which a rational function $\psi(x_1, \hdots, x_n)$ takes when operated on by all $n!$ substitutions are the roots of
an equation of degree $\rho$ whose coefficients are rational functions of the elementary symmetric functions.

\end{theorem}

{\bf Formation of Functions of a Given Group}\\

We define the Galois Function having $n!$ values for all the substiutionts of the symmetric group
\[ \psi_1 = \alpha_1 x_ 1 + \alpha_2 x_ 2 + \alpha_3 x_ 3 + \hdots \alpha_n x_ n,\] in which $\alpha_1, \alpha_2, \hdots, \alpha_n$ are $n$ distinct arbitrary constants. This function is
called the Galois Function. If we use all the $P$ substitutions of $H$ to form the functions $\psi_1, \psi_2, \hdots, \psi_P$ by the substitutions of $H$ anysymmetric funciton of these values will be unaltered by the substitutions of $H$. In particular $\phi_1 = (y-\psi_1)(y-\psi_2)\hdots(y-\psi_P)$ will be unaltered by the substutions of $H$, and altered to a different value by the substitutions not in $H$. The function $\phi_1$ is not therefore unaltered by the substitutions of any wider group that contains $H$ as a subgroup. Noting the forms of the expression of the sums of powers of the roots in terms of the coefficients, we see that instead of the coefficients of the powers of $y$ in $\phi_1$ we may take sums of powers of $\psi_1, \psi_2, \hdots, \psi_P$ up to the $P^{th}$ and deduce that one at least of such $P$ sums of powers will provide a function unlatered by the substitutions of $H$ and altered by any other substitution of the symmetric group.

\begin{theorem}
Every integral symmetric function of the distinct values of any integral multiple-valued function of $n$ elements is a symmetric function of the elements themselves.
\end{theorem}

\begin{proof}
Let $F(\phi_1, \phi_2, \hdots, \phi_\rho)$ be any rational integral symmetric function of the $\rho$-values. Any substitution $S$ whatever affecting the elements applied to these $\rho$-values
either leaves any function unchanged or replaces it by one of the others. No two of the resulting values can be equal, for if $\phi_i S$ were equal to $\phi_j S$ it would follow, by applying the substitution $S^{-1}$,  that $\phi_i = \phi_j$ which is contrary to the assumptions. Consequently the $\rho$ values of $\phi$ are reproduced by $S$ in some order or other. The symmetric function $F$ therefore remains unchanged by any substitutiont, and is consequently a symmetric element of the elements themselves. 
\end{proof}

\begin{corollary}
The $\rho$ distinct values of an integral multiple-valued function are roots of an equation whose coefficients are integral symmetric functions of the elements themselves.
\end{corollary}

\begin{theorem}[Lagrange Theorem on functions belonging to the same group]
Two rational functions belonging to the same group are rationally expressed each in terms of the other.
\label{FunctionsOfSameGroup}\end{theorem}

\begin{proof}
Let $\phi_1$ and $\psi_1$ be two functions belonging to the same group
\[ H = \{ I, S_2, S_3, \hdots, S_r \}, \] of order $r$ and degree $n$, each of these functions having $\rho$ distinct values, where $r*\rho = N$. Any substitution not contained in $H$
will convert $\phi_1$ into another of its values, say $\phi_k$ and at the same time convert $\psi_1$ into $\psi_k$. By operating all possible substitutions, $\rho$ pairs of values
$\phi_1, \psi_1 : \phi_2, \psi_2 : \hdots \phi_\rho, \psi_\rho$ are obtained. Now, in the first place, the rational function

\be \sum \phi^i \psi^j = \phi_1^i \psi_1^j + \phi_2^i \psi_2^j + \hdots + \phi_k^i \psi_k^j + \hdots + \phi_\rho^i \psi_\rho^j \label{psiphi}\ee
is clearly a symmetric function of the elements, for it appears by the same reasoning as before that any substitution acting on this function will merely permute the terms leaving the sum unaltered. Therefore $\sum \phi^i \psi^j$ is a symmetric function of the elements. \\

If we now take $j = 1$ and assign to $i$ all the values $0,1,2,\hdots,\rho-1$ in succession, we obtain the following $\rho$ equations linear in $\psi_1, \psi_2, \hdots, \psi_\rho$:
\be
\begin{array}{r c r c r c r r}
\psi_1 & + & \psi_2 & + & \hdots & + & \psi_\rho = & T_0 \\
\phi_1 \psi_1 & + & \phi_2\psi_2& + & \hdots & + &\phi_\rho \psi_\rho = & T_1 \\
\phi_1^2 \psi_1 & + & \phi_2^2\psi_2& + & \hdots & + &\phi_\rho^2 \psi_\rho = & T_2 \\
\hdots & \hdots & \hdots & \hdots & \hdots & \hdots &\hdots = & \hdots \\
\phi_1^{\rho-1} \psi_1 & + & \phi_2^{\rho-1}\psi_2& + & \hdots & + &\phi_\rho^{\rho-1} \psi_\rho = & T_{\rho-1} \\
\end{array}\label{linearSys}
\ee where $T_0, T_1, \hdots, T_{\rho-1}$ are all symmetric functions of $x_1, x_2, \hdots, x_n$.
\end{proof}
We can write the above as a linear system
\be
\begin{bmatrix}
1 & 1 & 1 & 1 & \hdots & 1 \\
\phi_1 & \phi_2 & \phi_3 & \phi_4 & \hdots & \phi_\rho \\
\phi_1^2 & \phi_2^2 & \phi_3^2 & \phi_4^2 & \hdots & \phi_\rho^2 \\
\vdots & \vdots & \vdots& \vdots & \vdots & \vdots \\
\phi_1^{\rho-1} & \phi_2^{\rho-1} & \phi_3^{\rho-1} & \phi_4^{\rho-1} & \hdots & \phi_\rho^{\rho-1} \\
\end{bmatrix}
\begin{bmatrix}
\psi_1 \\
\psi_2 \\
\vdots\\
\vdots\\
\psi_\rho
\end{bmatrix}
=
\begin{bmatrix}
T_0 \\
T_1 \\
\vdots\\
\vdots\\
T_{\rho-1}
\end{bmatrix}
\label{LagrangeSystem}\ee

Linear systems such as Eq.(\ref{LagrangeSystem}) can be solved using Cramer's rule. Alternatively, in order to understand better the relatonship between $\phi_i$ and $\psi_i$ let us consider
a method based on transforming non-homogenous linear systems into sets of corresponding homogeneous linear systems in each of the unknown parameters. To this effect consider 
the linear system \cite{Burnside}.
\be
\begin{array}{r r r r r r r r}
x               & + &             y   & + &                z     & = & T_0\\
\alpha x    & + &     \beta y  & + & \gamma z      & = & T_1 \\
\alpha^2 x & + & \beta^2 y  & + & \gamma^2 z  & = & T_2 \\
\end{array}
\label{ExampleLinearSystem}
\ee
By adding a fourth equation to the above we can isolate one of the variables, say $y$, and determine the functional dependence of $y$ on $\beta$ and the symmetric functions of the coefficients.
Adding $y = y$ to the above system and converting it into a system of homogenous equations by introducing a 4th unknown $w$ we obtain the system 

\be
\begin{array}{r r r r r r r r}
0               & + &             y   & + &                0     & = & y   \\
x               & + &             y   & + &                z     & = & T_0\\
\alpha x    & + &     \beta y  & + & \gamma z      & = & T_1 \\
\alpha^2 x & + & \beta^2 y  & + & \gamma^2 z  & = & T_2 \\
\end{array}
\ee

The above system can be converted into a Homogeneous System by moving the final column to the left-hand side of the equal sign. 

\be
\begin{array}{r r r r r r r r r r}
0               & + &             y   & + &                0     & - & y    & = & 0\\
x               & + &             y   & + &                z     & - & T_0& = & 0\\
\alpha x    & + &     \beta y  & + & \gamma z      & - & T_1 & = & 0\\
\alpha^2 x & + & \beta^2 y  & + & \gamma^2 z  & - & T_2 & = & 0\\
\end{array}
\ee

Writing the above system using a matrix of coefficients and introducing the 4th undetermined variable we find
\be
\begin{bmatrix}
0            & 1           & 0 & -y  \\
1            & 1           & 1 & -T_0  \\
\alpha    & \beta     & \gamma & -T_1  \\
\alpha^2 & \beta^2 & \gamma^2 & -T_2  \\
\end{bmatrix}
\begin{bmatrix}
x \\
y \\
z \\
w \\
\end{bmatrix}
=
\begin{bmatrix}
0 \\
0 \\
0 \\
0 \\
\end{bmatrix}
\label{BandP}\ee

In order for this homogenous system to be consistent with Eq.(\ref{ExampleLinearSystem}) the determinant of the matrix of coefficients must vanish. 
\begin{equation}
\begin{vmatrix}
0            & 1           & 0               & -y    \\
1            & 1           & 1               & -T_0\\
\alpha    & \beta     & \gamma    & -T_1\\
\alpha^2 & \beta^2 & \gamma^2 & -T_2\\
\end{vmatrix}\end{equation}
Since the determinant vanishes we can factor out the $-1$ from the last column and write the consistency condition as 

\begin{equation}
\begin{vmatrix}
0            & 1           & 0               &  y    \\
1            & 1           & 1               &  T_0\\
\alpha    & \beta     & \gamma    &  T_1\\
\alpha^2 & \beta^2 & \gamma^2 & T_2 \\
\end{vmatrix} = 0.
\end{equation}

Multiplying the above determinant by the non-singular determinant 
\begin{equation}
\begin{vmatrix}
1            & 1           & 1               &  0  \\
\alpha    & \beta     & \gamma    &  0   \\
\alpha^2 & \beta^2 & \gamma^2 & 0   \\
0            & 0           & 0               &  1   \\
\end{vmatrix} 
\end{equation}we find

\begin{equation}
\begin{vmatrix}
0            & 1           & 0               &  y     \\
1            & 1           & 1               &  T_0 \\
\alpha    & \beta     & \gamma    &  T_1 \\
\alpha^2 & \beta^2 & \gamma^2 & T_2  \\
\end{vmatrix} 
\begin{vmatrix}
1   & \alpha    & \alpha^2   &  0    \\
1   & \beta      & \beta^2    &  0    \\
1   & \gamma & \gamma^2 & 0    \\
0    & 0           & 0               &  1    \\
\end{vmatrix}
=
\begin{vmatrix}
1            & \beta     & \beta^2      &  y    \\
s_0        & s_1       & s_2             &  T_0\\
s_1        & s_2       & s_3             &  T_1\\
s_2        &s_3        & s_4             & T_2 \\
\end{vmatrix} = 0,
\label{BandPFinal}\end{equation}
where \[s_i = \alpha^i + \beta^i + \gamma^i\] and the $s_i$ are symmetric functions in $\alpha, \beta, \mbox{  and } \gamma$. \\ 

Since the $s_i$'s are sums of powers of $\alpha, \beta, \gamma$ and are symmetric functions, they can be expressed rationally in terms of the elementary symmetric functions. By expanding the final determinant in Eq. (\ref{BandPFinal}) along the top row and setting the expansion equal to zero to satisfy the consistency requirement we see that we have $y$ expressed as a quadratic function of $\beta$ along with the sums of powers of the three quantities $\alpha, \beta, \gamma$. Summarizing the results obtained -- we expand the above determinant along the top row and solve for $y$.

\be
y = {\begin{vmatrix}
s_1       & s_2             &  T_0\\
s_2       & s_3             &  T_1\\
s_3       & s_4             & T_2 \\
\end{vmatrix} -
\beta\begin{vmatrix}
s_0        & s_2             &  T_0\\
s_1        & s_3             &  T_1\\
s_2        & s_4             & T_2 \\
\end{vmatrix} 
+
\beta^2\begin{vmatrix}
s_0        & s_1       &  T_0\\
s_1        & s_2       &  T_1\\
s_2        &s_3        & T_2 \\
\end{vmatrix} 
\over 
\begin{vmatrix}
s_0        & s_1       & s_2\\
s_1        & s_2       & s_3\\
s_2        &s_3        & s_4\\
\end{vmatrix}}.
\label{yBeta}\ee
The right-hand-side depends quadratically on $\beta$ and other expressions which are symmetric functions of $\alpha, \beta, \mbox{  and  } \gamma$ and hence demonstrate the claim made in Theorem (\ref{FunctionsOfSameGroup}).\\

The denominator of Eq. (\ref{yBeta}) is a function of the squares of the differences of $\alpha, \beta, \mbox{  and  } \gamma$:  
\be\triangle(\alpha, \beta, \gamma) = 
\begin{vmatrix}
s_0        & s_1       & s_2\\
s_1        & s_2       & s_3\\
s_2        &s_3        & s_4\\
\end{vmatrix} = (\alpha - \beta)^2 (\alpha - \gamma)^2  (\beta - \gamma)^2
\ee Using this result we can write the solution for $y$ as
\be\triangle(\alpha, \beta, \gamma) y = \begin{vmatrix}
s_1       & s_2             &  T_0\\
s_2       & s_3             &  T_1\\
s_3       & s_4             & T_2 \\
\end{vmatrix} -
\beta\begin{vmatrix}
s_0        & s_2             &  T_0\\
s_1        & s_3             &  T_1\\
s_2        & s_4             & T_2 \\
\end{vmatrix} 
+
\beta^2\begin{vmatrix}
s_0        & s_1       &  T_0\\
s_1        & s_2       &  T_1\\
s_2        &s_3        & T_2 \\
\end{vmatrix} .
\ee\\


Returning to the original problem  Eq. (\ref{linearSys}), we can similarly express the solution for $\psi_j$ in terms of $\phi_j$ in the form of a determinant: 

\be
\begin{vmatrix}
1            & \phi_j     & \phi_j^2      &  \hdots & \phi_j^{\rho-1} & \psi_j  \\
s_0        & s_1        & s_2      &  \hdots & s_{\rho-1}              & T_0     \\
s_1        & s_2        & s_3      &  \hdots & s_\rho                     & T_1     \\
\hdots    & \hdots    & \hdots        &  \hdots & \hdots              & \hdots  \\
s_{\rho-1}            & s_\rho     & s_{\rho+2}      &  \hdots & s_{2\rho-2} & T_{\rho-1}   \\
\end{vmatrix} = 0,
\ee where $T_0, T_1, \hdots, T_{\rho-1}$ are all symmetric functions of $x_1, x_2, \hdots, x_n$ and $s_k = \phi_1^k + \phi_2^k + \phi_3^k + \hdots + \phi_\rho^k$.\\

The solution of these equations takes the form 
\be \triangle(\phi_1, \phi_2, \phi_3, \hdots, \phi_\rho) \psi_j = A_0 \phi_j^{\rho-1} + A_1 \phi_j^{\rho-2} + \hdots + A_{\rho-1},\ee
where $\triangle$ is the product of the square of the differences of all the $\phi$'s (the square of the Vandermonde determinant) and $A_0, A_1, \hdots, A_{\rho-1}$ are all symmetric in $x_1, x_2, \hdots, x_n$. \\

Therefore $\psi_j$ can be rationally expressed as a rational function of $\phi_j$ and the elementary symmetric functions of $x_1, x_2, \hdots, x_n$ (the coefficients of the polynomial that has $x_1, x_2, \hdots, x_n$ as roots). \\

{\bf Lagrange Interpolation Formula: } \\
Given a set of $k+1$ nodes $\{x_0, x_1, \hdots, x_k\}$, which must all be distinct, $x_j \ne x_m$ for indices $j\ne m$, the {\it Lagrange basis} for polynomials of degree $\le k$ for those nodes is the set of polynomials $\{ l_0(x), l_1(x), \hdots, l_k(x)\}$ each of degree $k$ which take values $l_j(x_m) = 0$ if $m\ne j$ and $l_j(x_j) = 1$. Using the Kronecker delta this can be written
$l_k(x_m) = \delta_{jm}$. Each basis polynomial can be explicitly described by the product: 

\begin{eqnarray*}
l_j(x) &=& {(x-x_0)\over (x_j-x_0)}{(x-x_1)\over (x_j-x_1)} \hdots {(x-x_{j-1})\over (x_j-x_{j-1})}{(x-x_{j+1})\over (x_j-x_{j+1})}\hdots{(x-x_k)\over (x_j-x_k)}\\
         &=& \prod_{\substack{0\le m\le k \\ m \ne j }}\,\, {x-x_m\over x_j - x_m}
\end{eqnarray*}

Notice that the numerator $\displaystyle \prod_{m\ne j} (x-x_m)$ has $k$ roots at the nodes $\{x_m\}_{m\ne k}$ while the denominator 
$\displaystyle \prod_{m \ne j} (x_j -x_m)$ scales the resulting polynomial so that $l_j(x_j) = 1$. \\

The Lagrange interpolating polynomial for those nodes through the corresponding values $\{ y_0, y_1, \hdots, y_k\}$ is the linear combination
\[ L(x) = \sum_{j=0}^k \, y_j l_j(x),\]
Each basis polynomial has degree $k$, so the sum $L(x)$ has degree $\le k$, and it interpolates the data because 
\[ L(x_m) = \sum_{j=0}^k y_j l_j(x_m) = \sum_{j=0}^k y_j \delta_{mj} = y_m,\]

\begin{theorem}[Extension of Lagrange's Theorem to Non Symmetric Functions] If two rational functions of any set of variables are such that one remains unchanged by all the substitutions of the group to
which the other belongs, the first is expressible by means of the second in the form of an integral polynomial whose coefficients are rational symmetric functions of the variables. \\

{\bf Another statement of Lagrange's Theorem: } \\

If a rational function $\phi(x_1,x_2, \hdots, x_n)$ remains unaltered by all the substitutions which leaves another rational function $\psi(x_1, x_2, \hdots, x_n)$ unaltered, then $\phi$ is
a rational function of $\psi$ and $c_1, c_2, \hdots, c_n$. Note that the group, $G$ to which $\phi$ belongs contains the group, $H$ to which $\psi$ belongs, i.e. $H\subseteq G$. $G$ is the
wider group. Also $\phi$ is therefore the {\it more symmetric} function -- a wider set of substitutions keep it unaltered.
\end{theorem}

\begin{proof}
\be
\begin{array}{llll | r | r}
I & h_2 & \hdots & h_P & \psi = \psi_1 & \phi = \phi_1 \\
g_2 & h_2 g_2 & \hdots & h_P g_2 & \psi_{g_2} = \psi_2 & \phi_{g_2} = \phi_2 \\
\hdots & \hdots & \hdots & \hdots & \hfill \hdots \hfill & \hfill \hdots \hfill\\
g_\nu & h_2 g_\nu & \hdots & h_P g_\nu & \psi_{g_\nu} = \psi_\nu & \phi_{g_\nu} = \phi_\nu \\
\end{array}\
\ee
The $\psi_1, \psi_2, \hdots, \psi_\nu$ are all distinct; but $\phi_1, \phi_2, ...\phi_nu$ need not be distinct since $\phi$ belongs to a group $G$ which may be bigger than $H$.
Under any substitution $s$ on $x_1, x_2, \hdots, x_n$, the functions $\psi_1, \psi_2, \hdots, \psi_\nu$ are merely permuted. Moreover, if $s$ replaces
$\psi_i$ by $\psi_j$, it replaces $\phi_i$ by $\phi_j$. Set (using the Lagrange Interpolation formula)

\begin{eqnarray*}
g(t) &=& (t-\psi_1)(t-\psi_2)\hdots, (t-\psi_\nu), \cr
\lambda(t) &=& g(t) \left( {\phi_1\over t- \psi_1} +{\phi_2\over t- \psi_2} + \hdots + {\phi_\nu\over t- \psi_\nu} \right)
\end{eqnarray*}so that $\lambda(t)$ is an integral function of degree $\nu-1$ in $t$. Since $\lambda(t)$ remains unaltered under every substitution of $s$, its coefficients
are rational symmetric functions of $x_1, x_2, \hdots, x_n$ and hence are rational functions of the elementary symmetric functions. \\

Now we're assuming here that $x_1, x_2, \hdots, x_n$ denote indeterminant quantities since $\psi_1, \psi_2, \hdots, \psi_\nu$ are algebraically distinct, so that $g'(\psi)$ is not idendically
zero. In the case where special values are assigned to $x_1, x_2, \hdots, x_n$ such that two or more of the functions $\psi_1, \psi_2, \hdots, \psi_\nu$ become numerically equal, then
$g'(\psi) = 0$, and $\phi$ is not a rational function of $\psi, c_1, ..., c_n$.
Assuming that we're working with indetermiant $x_1, x_2, \hdots, x_n$, and
taking $\psi_1 = \psi$ for $t$, we get
\[ \lambda(\psi_1) = \left[ (\psi_1 - \psi_2)(\psi_1-\psi_3)\hdots (\psi_1-\psi_\nu) \right]\phi_1\] and in general
\[ \phi = {\lambda(\psi) \over g'(\psi)}.\]
\end{proof}

\begin{corollary}
A function can always be found in terms of which any number of given functions can be rationally expressed.
\end{corollary}
The groups if the given functions have always one sub-group common to all, for the identity $S=1$ is common to all. Therefore the functions can all be expressed in terms of any one of the 
functions peculiar to the common sub-group. Therefore if $\phi, \psi, \xi, \hdots, $ are the given functions, then
\[ f = \alpha\phi + \beta\psi + \gamma\xi + \hdots,\] where $\alpha, \beta, \gamma, \hdots,$ are arbitrary constants, is one such function for the common subgroup ; for any substitution 
which lease it unaltered must leave $\phi, \psi, \xi, \hdots$, unaltered and must be common to the groups of $\phi_1, \psi_1, \xi_1, \hdots$. 

\begin{corollary}
Any rational function whatever can be rationally expressed in terms of a function having $n!$ distinct values ; in particular the Galois function. 
\end{corollary}
The Group of an $n!$-valued function, is included as a sub-group of every other permutation group.

\begin{corollary}
The roots themselves can be expressed rationally in terms of the Galois function.
\end{corollary}

\begin{proof}
Let $f(x) = 0$ be the equation whose roots are $x_1, x_2, \hdots, x_n$, all supposed unequal; and let $\psi_1$ be a known value of a a rational function $\psi$ of the roots which has $n!$ distinct values.\\

If all the roots except $x_1$ be permuted in every possible way, we obtain $\mu = (n-1)!$ distinct values of $\psi$ given by the equation
\[ F(x) = (\psi - \psi_1)(\psi - \psi_2)\hdots (\psi - \psi_\mu) = 0.\] 
The coefficients of this equation when expanded are symmetric functions of $x_2, x_3, \hdots, x_n$ and therefore can  be expressed in terms of the coefficients of 
\[ {f(x) \over x-x_1} = 0, \]
and will involve $x_1$ in a rational form along with the coefficients of $f(x)$ If the expanded equation be represented by $F(\psi, x_1) = 0$, 
we have $F(\psi_1, x_1) = 0$, since it is satisfied by $\psi = \psi_1$; also we have $f(x_1) = 0$, from which it follows that the equation $f(x) = 0$ and $F(\psi_1, x) = 0$ have a common root. 
It is easily seen that this is the only common root. If therefore we seek the greatest common divisor of $f(x)$ and $F(\psi_1, x)$, as all the remainders vanish for $x = x_1$ in particular the remainder of the first degree in 
$x$ of $f(x)$, in which the expression $\psi_1$ may be replaced by $\psi_2, \psi_3, \hdots \psi_\mu$ without altering its value.  
\end{proof}


\begin{corollary}
All the values of the Galois function can be expressed rationally in terms of any one of them. 
\end{corollary}

\begin{proof}
For all of these values belong to the same group, i.e., the identity.
\end{proof}
 